\section{Characters and Strings}

\subsection{Section Overview}
Text is a fundamental data type in programming. In this section, we'll explore how C++ handles characters and strings. We will cover the \texttt{char} type and associated functions, the traditional C-style strings (which are arrays of characters), and the modern, more powerful \texttt{std::string} class from the C++ Standard Library.

\subsection{Character Functions}
The \texttt{<cctype>} header provides a set of functions for working with characters. These functions can be used to test characters for certain properties (e.g., is it a letter, a digit, etc.) and to convert their case.
\begin{lstlisting}
	#include <cctype> // include this header

	// Some common functions:
	// isalpha(c): true if c is a letter
	// isdigit(c): true if c is a digit
	// isalnum(c): true if c is a letter or a digit
	// isspace(c): true if c is a whitespace character
	// toupper(c): returns uppercase version of c
	// tolower(c): returns lowercase version of c
\end{lstlisting}

\subsection{C-Style Strings}
In C and C++, a ``string'' can be represented as an array of characters. This is known as a \textbf{C-style string}. By convention, C-style strings are terminated by a special character, the null character (\texttt{\textbackslash{}0}), which has an ASCII value of 0.

\begin{lstlisting}
	char my_string[] {"Hello"}; // size is 6 (H,e,l,l,o,\0)
	char name[50]; // can hold up to 49 characters + null
\end{lstlisting}

\subsection{Working with C-style Strings}
The \texttt{<cstring>} header provides functions for working with C-style strings, such as copying, concatenating, and comparing them.

\begin{lstlisting}
	#include <cstring> // include this header

	char str1[50] {"Hello"};
	char str2[50] {" World"};
	char str3[50];

	strcpy(str3, str1); // copies str1 to str3
	strcat(str3, str2); // concatenates str2 onto str3

	std::cout << "str3: " << str3 << std::endl; // "Hello World"
	std::cout << "Length of str3: " << strlen(str3) << std::endl; // 11
\end{lstlisting}

\begin{remark}
	Many C-style string functions are unsafe because they do not perform bounds checking. \texttt{strcpy} and \texttt{strcat} can easily lead to buffer overflows if the destination array is not large enough. Modern C++ code should prefer \texttt{std::string}.
\end{remark}

\subsection{Coding Exercise 19: Using C-style Strings}
Declare two C-style strings. Use \texttt{strcpy} to copy the first string into a third string, and \texttt{strcat} to concatenate the second string onto the third. Finally, print the length of the final string.
\begin{lstlisting}
	#include <iostream>
	#include <cstring>

	int main() {
		char first_name[] = "John";
		char last_name[] = " Doe";
		char full_name[50];

		strcpy(full_name, first_name);
		strcat(full_name, last_name);

		std::cout << "Full name: " << full_name << std::endl;
		std::cout << "Length: " << strlen(full_name) << std::endl;

		return 0;
	}
\end{lstlisting}

\subsection{C++ Strings}
The C++ Standard Library provides the \texttt{std::string} class, which is a much safer and more convenient way to work with strings. To use it, you must include the \texttt{<string>} header.

\begin{lstlisting}
	#include <string>

	std::string s1 {"Hello"};
	std::string s2; // empty string
	std::string s3 {s1}; // copy of s1
	std::string s4 (s1, 1, 3); // "ell"
	std::string s5 (5, 'x'); // "xxxxx"
\end{lstlisting}

\subsection{Working with C++ Strings}
\texttt{std::string} objects can be concatenated with \texttt{+}, compared with relational operators, and accessed using array syntax \texttt{[]} or the \texttt{.at()} method. They also have many useful member functions.

\begin{lstlisting}
	#include <iostream>
	#include <string>

	int main() {
		std::string s1 {"Hello"};
		std::string s2 {"World"};

		std::string s3 = s1 + " " + s2; // "Hello World"
		std::cout << s3 << std::endl;

		if (s1 == "Hello") {
			std::cout << "s1 is Hello" << std::endl;
		}

		std::cout << "Length of s3: " << s3.length() << std::endl; // 11

		// Get input with spaces
		std::string full_name;
		std::cout << "Enter your full name: ";
		std::getline(std::cin, full_name);
		std::cout << "Your name is: " << full_name << std::endl;

		return 0;
	}
\end{lstlisting}

\subsection{Coding Exercise 20: Using C++ Strings --- Exercise 1}
Ask the user to enter their full name and age. Concatenate their name and age into a single \texttt{std::string} and print it.
\begin{lstlisting}
	#include <iostream>
	#include <string>

	int main() {
		std::string name;
		int age;
		std::cout << "Enter your full name: ";
		std::getline(std::cin, name);
		std::cout << "Enter your age: ";
		std::cin >> age;

		std::string result = name + " is "
			+ std::to_string(age) + " years old.";
		std::cout << result << std::endl;

		return 0;
	}
\end{lstlisting}

\subsection{Coding Exercise 21: Using C++ Strings --- Exercise 2}
Write a program that takes a \texttt{std::string} and replaces all occurrences of the character \texttt{'a'} with \texttt{'@'}.
\begin{lstlisting}
	#include <iostream>
	#include <string>

	int main() {
		std::string text = "A fast and accurate C++ program.";
		for (char &c : text) {
			if (c == 'a' || c == 'A') {
				c = '@';
			}
		}
		std::cout << text << std::endl;
		return 0;
	}
\end{lstlisting}

\subsection{Section Challenge}
Write a program that implements a simple substitution cipher. Ask the user for a secret message (a string). Create two strings representing the alphabet and a substitution key. Encrypt the message by substituting characters and print it. Then, decrypt the encrypted message and print it to verify.

\subsection{Section Challenge --- Solution}
\begin{lstlisting}
	#include <iostream>
	#include <string>

	int main() {
		std::string alphabet {
			"abcdefghijklmnopqrstuvwxyz"
			"ABCDEFGHIJKLMNOPQRSTUVWXYZ"};
		std::string key {
			"XZNLWEBGJHQDYVTKFUOMPCIASRxznlwebgjhqdyvtkfuompciasr"};

		std::string secret_message;
		std::cout << "Enter your secret message: ";
		std::getline(std::cin, secret_message);

		std::string encrypted_message;
		for (char c : secret_message) {
			size_t position = alphabet.find(c);
			if (position != std::string::npos) {
				encrypted_message += key[position];
			} else {
				encrypted_message += c;
			}
		}
		std::cout << "\nEncrypting message..." << std::endl;
		std::cout << "Encrypted message: "
			<< encrypted_message << std::endl;

		std::string decrypted_message;
		for (char c : encrypted_message) {
			size_t position = key.find(c);
			if (position != std::string::npos) {
				decrypted_message += alphabet[position];
			} else {
				decrypted_message += c;
			}
		}
		std::cout << "\nDecrypting message..." << std::endl;
		std::cout << "Decrypted message: "
			<< decrypted_message << std::endl;

		return 0;
	}
\end{lstlisting}

\subsection{Quiz 7: Section 10 Quiz}
\begin{enumerate}
	\item What character terminates a C-style string?
	\begin{enumerate}
		\item[a)] \texttt{\textbackslash{}n} (newline)
		\item[b)] \texttt{\textbackslash{}t} (tab)
		\item[c)] \texttt{\textbackslash{}0} (null character)
		\item[d)] \texttt{.} (period)
	\end{enumerate}

	\item Which header file is required to use the \texttt{std::string} class?
	\begin{enumerate}
		\item[a)] \texttt{<cstring>}
		\item[b)] \texttt{<string>}
		\item[c)] \texttt{<iostream>}
		\item[d)] \texttt{<cctype>}
	\end{enumerate}

	\item Which of the following is a major advantage of \texttt{std::string} over C-style strings?
	\begin{enumerate}
		\item[a)] Automatic memory management and bounds checking on some methods.
		\item[b)] They are always faster.
		\item[c)] They use less memory.
		\item[d)] They are compatible with C functions.
	\end{enumerate}
\end{enumerate}

\textbf{Answers:} 1.\ (c), 2.\ (b), 3.\ (a)

\subsection{Assignment 1: Challenge Assignment --- Letter Pyramid}
Write a program that asks the user to enter a string of characters. The program should then display a pyramid of letters from the string. For example, if the user enters ``ABCDE'', the program should display:
\begin{center}
\begin{verbatim}
    A
   ABA
  ABCBA
 ABCDCBA
ABCDEDCBA
\end{verbatim}
\end{center}
\begin{lstlisting}
	#include <iostream>
	#include <string>

	int main() {
		std::string letters;
		std::cout << "Enter a string of letters: ";
		std::cin >> letters;

		size_t num_letters = letters.length();
		for (size_t i = 0; i < num_letters; ++i) {
			// Print leading spaces
			for (size_t j = 0; j < num_letters - i - 1; ++j) {
				std::cout << " ";
			}
			// Print increasing part
			for (size_t j = 0; j <= i; ++j) {
				std::cout << letters[j];
			}
			// Print decreasing part
			for (size_t j = i; j > 0; --j) {
				std::cout << letters[j - 1];
			}
			std::cout << std::endl;
		}

		return 0;
	}
\end{lstlisting}
